\documentclass[10pt,openany,a4paper]{article}
\usepackage{tikz, pgfplots, tkz-euclide,calc}
    \usetikzlibrary{patterns,snakes,shapes.arrows}
\usepackage{fancyhdr}
\usepackage{enumerate,hyperref,enumitem}
\usepackage{cancel}
\usepackage{varwidth}

% TAMBAHKAN PACKAGE SENDIRI KALAU KURANG

\usepackage{geometry}
\geometry{
	total = {210mm, 148.5mm},
	left = 25mm,
	right = 25mm,
	top = 25mm,
	bottom = 25mm,
}

% \fancyhead{}
% \fancyfoot{}
% \fancyhead[r]{}
% \fancyhead[l]{\fbox{\large{\textbf{@ahmadzakiyudin\_}}}}
\renewcommand{\headrulewidth}{0pt}
\renewcommand{\footrulewidth}{0pt}
\usepackage{graphicx}
\usepackage{subcaption}
\usepackage{soul}
\usepackage{hyperref}
\usepackage{float}
\usepackage{setspace}
\usepackage{afterpage}
\usepackage{xcolor}
\usepackage[nosectionbib]{apacite}
\usepackage{lmodern}
\usepackage{mathtools}
\usepackage{upgreek}
\usepackage{siunitx}
\usepackage{physics}
\usepackage{hyperref}
\renewcommand{\baselinestretch}{1.25}
\usepackage{multicol}
\usepackage{ragged2e,scalerel,amsmath,multicol,gensymb,setspace,
fancyhdr,amsfonts,tikz,pgfplots,nccmath,enumerate,verbatim,amsthm,physics}
\usepgfplotslibrary{polar,fillbetween}
\usepgflibrary{shapes.geometric}
\usetikzlibrary{calc,patterns,arrows}
\newcommand\mylog[1]{\mathop{{}^{#1}\mathrm{log}}}
\pgfplotsset{compat=1.15}
\pgfplotsset{my style/.append style={axis x line=middle, axis y line=
middle, xlabel={$x$}, ylabel={$y$}, axis equal }}
\usepackage{etoolbox}
\makeatother
\newtheorem{defn}{Definisi}[section]
\newtheorem{teo}[defn]{Teorema}
\newtheorem{thm}{Teorema}[section]
\newtheorem{lemma}{Lemma}
\newtheorem{lemmas}[defn]{Lemma}

\newtheorem{cor}[defn]{Akibat}
\newtheorem{asumsi}[defn]{Asumsi}
\theoremstyle{definition}
\newtheorem{con}{Contoh}[section]
\theoremstyle{definition}
\theoremstyle{plain}
\newtheorem{prop}{Proposisi}[section]
\newcommand{\longdiv}{\smash{\mkern-0.43mu\vstretch{1.31}{\hstretch{.7}{)}}\mkern-5.2mu\vstretch{1.31}{\hstretch{.7}{)}}}}

\pagestyle{fancy}
\fancyhf{}
% \lhead{Halaman \thepage}
\cfoot{\href{https://instagram.com/ahmadzakiyudin_/}{@ahmadzakiyudin\_}}
\hypersetup{
    colorlinks=true,
    linkcolor=blue,
    filecolor=blue,      
    urlcolor=blue,
}
\newenvironment{jawab}[2][tes]{
\fbox{\begin{minipage}{14.65cm}
\textbf{Kata kunci materi:} #1.\\
\textbf{Pembahasan:}\\ #2
\end{minipage}}}{}

\newenvironment{jawab2}[2][tes]{
\fbox{\begin{minipage}{14.65cm}
 #2
\end{minipage}}}{}

\newenvironment{jawaban}[2][tes]{
\colorlet{oldcolor}{.}
        \color{red}
        \setlength{\fboxrule}{1pt}\boxed{\color{oldcolor}#2}\color{oldcolor}.
}{}

\newenvironment{jawaban2}[2][tes]{
\colorlet{oldcolor}{.}
        \color{red}
        \setlength{\fboxrule}{1pt}\fbox{\color{oldcolor}#2}\color{oldcolor}.
}{}

\newenvironment{penting}[2][tes]{
\colorlet{oldcolor}{.}
\color{blue}
\fbox{\color{oldcolor}#2}\color{oldcolor}
}{}
\renewcommand{\proofname}{Bukti}
\renewcommand{\thethm}{\arabic{chapter}.\arabic{thm}}

\newcommand{\normlp}[1]{\left\|#1\right\|_{L^p_{\lambda}(0,T;H)}}% Fungsi norm (||x||)
\newcommand{\R}{\mathbb{R}}
\fancyhead[l]{\fbox{\large{\textbf{PAKET A}}}}
\renewcommand{\natural}{\mathbb{N}}
\newcommand{\Xpl}{X^R_{p,\lambda}}
\newcommand{\Xpls}{X^R_{p,\lambda^*}}
\newcommand{\qpl}{q^R_{p,\lambda}}
\newcommand{\qpls}{q^R_{p,\lambda^*}}

\begin{document}

    \begin{center}
	{\underline{\textbf{\MakeUppercase{Try Out KUIS 1 Kalkulus 1}}}}
    \end{center}

    \begin{center}
	\begin{tabular}{lcl}
		Mata kuliah/SKS & : & Kalkulus 1 ( SM234201 ) / 3 SKS\\
		Hari, Tanggal & : & Sabtu, 19 September 2026\\
		Waktu & : & (100 menit)\\
		Sifat & : & Sesuai Peraturan Kuis\\
        Dosen & : & Pak Nurul\\
		% Kelas & : & 1-13, 101\\
  ~
	\end{tabular}
    
    \noindent\rule{\textwidth}{2.pt}
        \end{center}
\begin{enumerate}
    \item Tentukan semua nilai $x$ yang memenuhi ketaksamaan berikut
    \begin{enumerate}
    \item $\dfrac{1}{x+2}\geq \dfrac{3}{x-3}$.
    \item $2\leq |3x+7|<15$.
    \end{enumerate}
    \item Diberikan titik $A(1,2)$ dan $B(3,4)$. Tentukan persamaan garis yang melalui titik $A$ dan tegak lurus dengan garis yang melalui titik $B$ dan $C(5,6)$.
    \item Gambarkan grafik dari fungsi $f(x)=\begin{cases}
    (x+2)^2, & \text{jika } x<-1\\
    \sqrt{x+1}+1, & \text{jika } x\geq -1.
    \end{cases}$
    \item Diberikan fungsi $f(x)=\sqrt{2x-3}$ dan $g(x)=2x+3$ didefinisikan pada interval $-3\leq x\leq 5$. Dapatkan $(f\circ g)(x)$ dan tentukan domainnya.
\end{enumerate}
\newpage 
\fancyhead{l}{\fbox{\large{\textbf{PAKET B}}}}
 \begin{center}
    {\underline{\textbf{\MakeUppercase{Try Out KUIS 1 KALKULUS 1}}}}
    \end{center}
    \begin{center}
	\begin{tabular}{lcl}
		Mata kuliah/SKS & : & Kalkulus 2 ( SM234201 ) / 3 SKS\\
		Hari, Tanggal & : & Selasa, 31 Maret 2026\\
		Waktu & : & (100 menit)\\
		Sifat & : & Sesuai Peraturan Kuis\\
        Dosen & : & Pak Nurul\\
		% Kelas & : & 1-13, 101\\
  ~
	\end{tabular}
    
    \noindent\rule{\textwidth}{2.pt}
        \end{center}
        \begin{enumerate}
            \item Tentukan semua nilai $x$ yang memenuhi ketaksamaan berikut
            \begin{enumerate}
                \item $\dfrac{1}{x+2}\geq 3$.
                \item $1\leq |2x-8|<10$.
            \end{enumerate}
            \item Diberikan persamaan lingkaran $x^2+y^2-8x-6y+9=0$. Tentukan koordinat titik pusat dan panjang jari-jari lingkaran tersebut. Gambarkan lingkaran tersebut pada bidang koordinat.
            \item Gambarkan grafik dari fungsi $f(x)=\begin{cases}
    -x^2+4, & \text{jika } x<0\\
    \sqrt{x}+4, & \text{jika } x\geq 0.
    \end{cases}$
            \item Dapatkan nilai $k$ sehingga fungsi $f(x)=x^2+4x-1$ memiliki invers pada interval $[k,\infty)$. Tentukan juga invers dari fungsi tersebut beserta domainnya.
        \end{enumerate}
%         \begin{enumerate}
%             \item Dapatkan $\dfrac{dy}{dx}$ dari $y=\dfrac{e^{-x}+\ln x}{\cosh (5x)}$.
%             \item Dapatkan $\dfrac{dy}{dx}$ dari $y=e^{x\tan^{-1}x}$.
%             \item Hitung 
%             \begin{align*}
%                 \int_0^{\frac{1}{2}\ln 3} \frac{e^x}{1+e^{2x}}\, dx.
%             \end{align*}
%             \item Hitung 
%             \begin{align*}
%                 \int \cosh (4x)\sqrt{2+7\sinh(4x)}\, dx
%             \end{align*}
%         \end{enumerate}

% \newpage
% \fancyhead[l]{\fbox{\large{\textbf{PAKET C}}}}
%  \begin{center}
% 	{\underline{\textbf{\MakeUppercase{Try Out KUIS 1 KALKULUS 2}}}}
%     \end{center}

%     \begin{center}
% 	\begin{tabular}{lcl}
% 		Mata kuliah/SKS & : & Kalkulus 2 ( SM234201 ) / 3 SKS\\
% 		Hari, Tanggal & : & Selasa, 31 Maret 2026\\
% 		Waktu & : & (100 menit)\\
% 		Sifat & : & Terserah\\
% 		% Kelas & : & 1-13, 101\\
%   ~
% 	\end{tabular}
    
%     \noindent\rule{\textwidth}{2.pt}
%         \end{center}
%         \begin{enumerate}
%              \item  Dapatkan $\dfrac{dy}{dx}$ dari $y=\dfrac{\sin(2x)\sqrt{3x+\ln x}}{(x^3+3)^5}$.
%             \item Hitung 
%     \begin{align*}
%         \int \dfrac{x}{9+4x^4}\, dx.
%     \end{align*}
%             \item Hitung 
%     \begin{align*}
%         \int \frac{\sin \theta }{\sqrt{\cos^2\theta -1}}d \theta.
%     \end{align*}
%             \item Hitung 
%             \begin{align*}
%                 \int t \tan^{-1} t \, dt.
%             \end{align*}
%             \item Hitung
%             \begin{align*}
%                 \int \dfrac{t^2+5}{t^2-4t+5}\, dt
%             \end{align*}
%         \end{enumerate}

% \newpage
% \fancyhead[l]{\fbox{\large{\textbf{PAKET D}}}}
%  \begin{center}
% 	{\underline{\textbf{\MakeUppercase{Try Out KUIS 1 KALKULUS 2}}}}
%     \end{center}

%     \begin{center}
% 	\begin{tabular}{lcl}
% 		Mata kuliah/SKS & : & Kalkulus 2 ( SM234201 ) / 3 SKS\\
% 		Hari, Tanggal & : & Selasa, 31 Maret 2026\\
% 		Waktu & : & (100 menit)\\
% 		Sifat & : & Terserah\\
% 		% Kelas & : & 1-13, 101\\
%   ~
% 	\end{tabular}
    
%     \noindent\rule{\textwidth}{2.pt}
%         \end{center}
%         \begin{enumerate}
%              \item  Dapatkan $\dfrac{dy}{dx}$ dari $y=(x^2+2)^{2x+3\cosh x}$.
%             \item Hitung 
%     \begin{align*}
%         \int \dfrac{dt}{\sqrt{3e^{2t}-1}}
%     \end{align*}
%             \item Hitung 
%     \begin{align*}
%         \int \dfrac{dx}{x\sqrt{4(\ln x)^2-1}}.
%     \end{align*}
%             \item Hitung 
%             \begin{align*}
%                 \int \ln(2t^2+3) \, dt.
%             \end{align*}
%             \item Hitung
%             \begin{align*}
%                 \int \dfrac{t^2}{(t+2)^3}\, dt
%             \end{align*}
%         \end{enumerate}
\end{document}
